% GENERATED FILE. Edit canonical lessons, then run npm run generate:books.
% canonical-source-hash: fnv1a64:6161e050eaea4ad2
% canonical-lessons: KA-C249-purva, KA-C249-pascima, KA-C249-daksina, KA-C249-eduru, KA-C249-jille

\chapter{East, West, South, Opposite, District}
\label{ch:ka-east-west-south-opposite-district}
\hlchaptermodality{249}

\begin{chapteropening}
\textbf{By the end of this chapter:} \emph{I can say east, west, south, opposite and a district.}

Complete the last lesson of chapter 249: I can say east, west, south, opposite and a district.
\end{chapteropening}

\section[\texorpdfstring{\kn{ಪೂರ್ವ} (pūrva)}{ಪೂರ್ವ (pūrva)}]{\kn{ಪೂರ್ವ} (pūrva) — east}

\label{lesson:KA-C249-purva}



\begin{quote}
Before the new one: say the Kannada for an auspicious time, then the Kannada for one and a half.
\end{quote}

\subsection*{You'll want to know: \kn{ಪೂರ್ವ}}
\textbf{\kn{ಪೂರ್ವ}} — \emph{pūrva} — \textquotedblleft{}east\textquotedblright{}.

\begin{grammarlens}[title={What you've built}]
One more word for this chapter.
\end{grammarlens}

\subsection*{Your turn}
\begin{itemize}
\raggedright
  \item \emph{Say it:} \emph{pūrva}
  \item \emph{Say it:} \emph{pūrva}, once more
  \item \emph{Say it:} say \emph{pūrva}
  \item \emph{Recall:} say the Kannada for an auspicious time, then the Kannada for one and a half, then say \emph{pūrva} again
\end{itemize}

\begin{tcolorbox}[breakable,colback=teal!4,colframe=teal!35!black,title={Before you move on}]
What does \kn{ಪೂರ್ವ} mean? (\textquotedblleft{}East\textquotedblright{}.) Say it once more.
\end{tcolorbox}
\section[\texorpdfstring{\kn{ಪಶ್ಚಿಮ} (paścima)}{ಪಶ್ಚಿಮ (paścima)}]{\kn{ಪಶ್ಚಿಮ} (paścima) — west}

\label{lesson:KA-C249-pascima}



\begin{quote}
Before the new one: say the Kannada for one and a half, then the Kannada for east.
\end{quote}

\subsection*{You'll want to know: \kn{ಪಶ್ಚಿಮ}}
\textbf{\kn{ಪಶ್ಚಿಮ}} — \emph{paścima} — \textquotedblleft{}west\textquotedblright{}.

\begin{grammarlens}[title={What you've built}]
One more word for this chapter.
\end{grammarlens}

\subsection*{Your turn}
\begin{itemize}
\raggedright
  \item \emph{Say it:} \emph{paścima}
  \item \emph{Say it:} \emph{paścima}, once more
  \item \emph{Say it:} say \emph{paścima}
  \item \emph{Recall:} say the Kannada for one and a half, then the Kannada for east, then say \emph{paścima} again
\end{itemize}

\begin{tcolorbox}[breakable,colback=teal!4,colframe=teal!35!black,title={Before you move on}]
What does \kn{ಪಶ್ಚಿಮ} mean? (\textquotedblleft{}West\textquotedblright{}.) Say it once more.
\end{tcolorbox}
\section[\texorpdfstring{\kn{ದಕ್ಷಿಣ} (dakṣiṇa)}{ದಕ್ಷಿಣ (dakṣiṇa)}]{\kn{ದಕ್ಷಿಣ} (dakṣiṇa) — south}

\label{lesson:KA-C249-daksina}



\begin{quote}
Before the new one: say the Kannada for east, then the Kannada for west.
\end{quote}

\subsection*{You'll want to know: \kn{ದಕ್ಷಿಣ}}
\textbf{\kn{ದಕ್ಷಿಣ}} — \emph{dakṣiṇa} — \textquotedblleft{}south\textquotedblright{}.

\begin{grammarlens}[title={What you've built}]
One more word for this chapter.
\end{grammarlens}

\subsection*{Your turn}
\begin{itemize}
\raggedright
  \item \emph{Say it:} \emph{dakṣiṇa}
  \item \emph{Say it:} \emph{dakṣiṇa}, once more
  \item \emph{Say it:} say \emph{dakṣiṇa}
  \item \emph{Recall:} say the Kannada for east, then the Kannada for west, then say \emph{dakṣiṇa} again
\end{itemize}

\begin{tcolorbox}[breakable,colback=teal!4,colframe=teal!35!black,title={Before you move on}]
What does \kn{ದಕ್ಷಿಣ} mean? (\textquotedblleft{}South\textquotedblright{}.) Say it once more.
\end{tcolorbox}
\section[\texorpdfstring{\kn{ಎದುರು} (eduru)}{ಎದುರು (eduru)}]{\kn{ಎದುರು} (eduru) — opposite, facing}

\label{lesson:KA-C249-eduru}



\begin{quote}
Before the new one: say the Kannada for west, then the Kannada for south.
\end{quote}

\subsection*{You'll want to know: \kn{ಎದುರು}}
\textbf{\kn{ಎದುರು}} — \emph{eduru} — \textquotedblleft{}opposite, facing\textquotedblright{}.

\begin{grammarlens}[title={What you've built}]
One more word for this chapter.
\end{grammarlens}

\subsection*{Your turn}
\begin{itemize}
\raggedright
  \item \emph{Say it:} \emph{eduru}
  \item \emph{Say it:} \emph{eduru}, once more
  \item \emph{Say it:} say \emph{eduru}
  \item \emph{Recall:} say the Kannada for west, then the Kannada for south, then say \emph{eduru} again
\end{itemize}

\begin{tcolorbox}[breakable,colback=teal!4,colframe=teal!35!black,title={Before you move on}]
What does \kn{ಎದುರು} mean? (\textquotedblleft{}Opposite, facing\textquotedblright{}.) Say it once more.
\end{tcolorbox}
\section[\texorpdfstring{\kn{ಜಿಲ್ಲೆ} (jille)}{ಜಿಲ್ಲೆ (jille)}]{\kn{ಜಿಲ್ಲೆ} (jille) — a district}

\label{lesson:KA-C249-jille}



\begin{quote}
Before the new one: say the Kannada for south, then the Kannada for opposite.
\end{quote}

\subsection*{You'll want to know: \kn{ಜಿಲ್ಲೆ}}
\textbf{\kn{ಜಿಲ್ಲೆ}} — \emph{jille} — \textquotedblleft{}a district\textquotedblright{}.

\begin{grammarlens}[title={What you've built}]
One more word for this chapter.
\end{grammarlens}

\subsection*{Your turn}
\begin{itemize}
\raggedright
  \item \emph{Say it:} \emph{jille}
  \item \emph{Say it:} \emph{jille}, once more
  \item \emph{Say it:} say \emph{jille}
  \item \emph{Recall:} say the Kannada for south, then the Kannada for opposite, then say \emph{jille} again
\end{itemize}

\begin{tcolorbox}[breakable,colback=teal!4,colframe=teal!35!black,title={Before you move on}]
What does \kn{ಜಿಲ್ಲೆ} mean? (\textquotedblleft{}A district\textquotedblright{}.) Say it once more.
\end{tcolorbox}
